\subsection{Conclusiones}
\begin{itemize}
    \item La arquitectura modular que usamos en el proyecto nos permitió dividir todo bien claro. La lógica del negocio va por un lado, el manejo de los datos por otro y la interfaz queda separada. Gracias a eso mantener el sistema resulta mucho más fácil, se cometen menos errores porque las partes no se interfieren entre sí y si más adelante hay que crecer o agregar funcionalidades nuevas se puede hacer sin tocar ni romper lo que ya está funcionando.
    
    \item El uso de Firebase Realtime Database combinado con el patrón Repository nos permitió sincronizar de forma rápida y sin fallos los pedidos, sus estados y las ubicaciones. Gracias a eso tanto los administradores como los repartidores siempre trabajan con la información más reciente y actualizada en tiempo real.

    \item La integración del sistema de geolocalización y monitoreo en vivo nos ayudó muchísimo a mejorar el control de toda la logística. Ahora se puede ver en tiempo real las rutas que siguen los repartidores, el estado actual de cada entrega y exactamente dónde está cada uno. Esto hace que sea mucho más fácil tomar decisiones operativas rápidas y acertadas para que todo funcione mejor.
\end{itemize}

\subsection{Recomendaciones}
\begin{itemize}
    \item Hay que seguir mejorando la arquitectura modular del proyecto. Para eso, lo ideal es ir agregando pruebas automatizadas y también una buena documentación técnica. De esta forma se mantiene todo más estable, se evitan problemas inesperados y cuando toque hacer mejoras o agregar cosas nuevas en el futuro, va a ser mucho más fácil y rápido.

    \item Se recomienda optimizar cómo se maneja la data en tiempo real. Para eso, lo mejor es monitorear bien el rendimiento, controlar que las sincronizaciones vayan fluidas y ver si se pueden mejorar las actualizaciones para que sean más eficientes. Así se evita que la información se vuelva inconsistente entre los dispositivos.
    
    \item Se sugiere ampliar las partes de análisis y visualización de la app, agregando más KPIs, reportes detallados y métricas relacionadas con la logística. De esta forma se ayuda a planificar mejor las operaciones y los administradores pueden tomar decisiones con información más clara y completa, en lugar de ir a ciegas.
\end{itemize}